\section*{Exercises}
\addcontentsline{toc}{section}{Exercises}
\subsection*{A. Theory}
\begin{exercise}\exerciseid{19.A1} Explain why a low-rank signal plus full-rank noise may still have a visible singular-value elbow. \end{exercise}
\subsection*{B. By Hand}
\begin{exercise}\exerciseid{19.B1} Suppose a matrix has singular values $10,7,1,0.8,0.7$. Which rank would you try first for denoising, and why? \end{exercise}
\subsection*{C. Python / NumPy}
\begin{exercise}\exerciseid{19.C1} Generate a rank-$2$ matrix plus Gaussian noise. Compute singular values and compare rank-$1$, rank-$2$, and rank-$5$ reconstructions. \end{exercise}
\subsection*{D. Challenge}
\begin{exercise}\exerciseid{19.D1} Explain why too small a truncation rank underfits while too large a rank leaves noise in the reconstruction. \end{exercise}
